% 按用户指南复制三种版式，修改封面、层级、图片与两栏比例。
\documentclass[aspectratio=169,11pt,fontset=none]{ctexbeamer}
\usetheme{hitacademic}
\hitnosectionpages
\hitoutlinesubs
\hitfigures{example/images/}
\title[作者流程]{作者操作流程示例}
\author{示例作者}
\institute{示例单位}
\date{}
\begin{document}
\TitleSlide
\OutlineSlide[title={汇报提纲}]

\section{问题背景}
\subsection{问题范围}
\begin{frame}{输入与目标}
  \subhead{研究对象}
  \begin{itemize}
    \item 比较定义域 $[0,1]$ 上的两个函数。
    \item 说明线性关系与平方关系的差异。
  \end{itemize}
  \insight{比较采用相同的输入范围。}
\end{frame}

\section{函数关系}
\subsection{线性函数}
\begin{frame}{横图与说明}
  \hitcolumnratio{.6}
  \begin{columns}[T,onlytextwidth]
    \begin{column}{\hitleftwidth}
      \subhead{定义}
      线性函数满足 $y=x$。
    \end{column}
    \begin{column}{\hitrightwidth}
      \fig{linear.png}
      \captiontext{来源：模板解析函数图。}
    \end{column}
  \end{columns}
\end{frame}

\section{比较与讨论}
\subsection{平方函数}
\begin{frame}{竖图与横图对比}
  \hitcolumnratio{.45}
  \begin{columns}[T,onlytextwidth]
    \begin{column}{\hitleftwidth}
      \fig[height=.43\textheight]{quadratic.jpg}
      \captiontext{平方关系，来源：模板解析函数图。}
    \end{column}
    \begin{column}{\hitrightwidth}
      \fig[height=.43\textheight]{linear.pdf}
      \captiontext{线性关系，来源：模板解析函数图。}
    \end{column}
  \end{columns}
  \insight{区间内部有 $x^2<x$，端点处相等。}
\end{frame}
\EndSlide{感谢聆听}
\end{document}
